% Telos shared preamble.
%
% Every Telos publication is designed to be printed on a home laser printer,
% one sheet at a time, and carried into a boat or a garage. Two constraints
% follow from that and govern everything below:
%
%   1. Pure monochrome. No value is ever a hue. Emphasis is carried by rule
%      weight, fill, and type, never by color, so a grayscale or black-only
%      printer loses nothing.
%   2. Card density. A "one-pager" is a hard contract: the leaf must fit on a
%      single side of US Letter. The layout is therefore tight by default and
%      the environments below are sized to leave room for a plate.

\documentclass[10pt,letterpaper]{article}

\usepackage[margin=0.55in,includefoot,footskip=0.24in]{geometry}
\usepackage[T1]{fontenc}
\usepackage[utf8]{inputenc}
\usepackage{lmodern}
\usepackage{microtype}
\usepackage{array}
\usepackage{booktabs}
\usepackage{longtable}
\usepackage{tabularx}
\usepackage{enumitem}
\usepackage{needspace}
\usepackage{multicol}
\usepackage{ragged2e}
\usepackage{graphicx}
\usepackage{wrapfig}
\usepackage{xcolor}
\usepackage{fancyhdr}
\usepackage{titlesec}
\usepackage{hyperref}
\usepackage[most]{tcolorbox}
\usepackage{tikz}
\usetikzlibrary{arrows.meta,calc,positioning,shapes.geometric,decorations.pathmorphing,patterns,fadings}

% Semantic names are kept so a document reads intelligibly, but every value is
% black, white, or a neutral gray that survives a black-only print driver.
\definecolor{telosink}{gray}{0}
\definecolor{telospaper}{gray}{1}
\definecolor{telosrule}{gray}{0.35}
\definecolor{telosfill}{gray}{0.90}
\definecolor{telosdeep}{gray}{0.72}
\definecolor{telosmute}{gray}{0.45}

\hypersetup{hidelinks,pdfborder={0 0 0}}

% Keep an unchanged source byte-reproducible across rebuilds: no build-clock
% creation date and no random trailer ID.
\pdfinfoomitdate=1
\pdftrailerid{}

\setlength{\parindent}{0pt}
\setlength{\parskip}{0.42em}
\setlength{\tabcolsep}{4pt}
\setlength{\columnsep}{1.1em}
\renewcommand{\arraystretch}{1.18}
\setlist[itemize]{leftmargin=1.15em,itemsep=0.12em,topsep=0.2em,parsep=0pt}
\setlist[enumerate]{leftmargin=1.5em,itemsep=0.18em,topsep=0.2em,parsep=0pt}
\setlist[description]{leftmargin=0pt,itemsep=0.16em,topsep=0.2em,parsep=0pt,
  font=\normalfont\bfseries}

% ---------------------------------------------------------------------------
% Document identity. Each leaf declares these before \begin{document} so the
% running head, the colophon, and the PDF metadata all agree.
% ---------------------------------------------------------------------------
\newcommand{\TelosProject}{Telos}
\newcommand{\TelosDocTitle}{}
\newcommand{\TelosDocSubtitle}{}
\newcommand{\TelosRevision}{}

\newcommand{\telosmeta}[4]{%
  \renewcommand{\TelosProject}{#1}%
  \renewcommand{\TelosDocTitle}{#2}%
  \renewcommand{\TelosDocSubtitle}{#3}%
  \renewcommand{\TelosRevision}{#4}%
  \hypersetup{pdftitle={#2},pdfsubject={#3; version #4},
    pdfkeywords={Telos, version #4},pdfauthor={Telos}}%
}

\pagestyle{fancy}
\fancyhf{}
\renewcommand{\headrulewidth}{0pt}
\renewcommand{\footrulewidth}{0.4pt}
\fancyfoot[L]{\footnotesize\textsc{\TelosProject}}
\fancyfoot[C]{\footnotesize\TelosDocTitle}
\fancyfoot[R]{\footnotesize\thepage}

% The first page of a one-pager carries the masthead instead of a running foot,
% so its footer is the compact provenance line.
\fancypagestyle{telosfirst}{%
  \fancyhf{}%
  \renewcommand{\headrulewidth}{0pt}%
  \renewcommand{\footrulewidth}{0.4pt}%
  \fancyfoot[L]{\footnotesize\textsc{\TelosProject}}%
  \fancyfoot[C]{\footnotesize\TelosRevision}%
  \fancyfoot[R]{\footnotesize\thepage}%
}

% ---------------------------------------------------------------------------
% Masthead. A heavy rule, the title, a light subtitle, a thin rule. The
% optional argument is a right-aligned tag (a species' scientific name, a
% lesson number) that sits on the title's baseline.
% ---------------------------------------------------------------------------
\newcommand{\telosmasthead}[1][]{%
  \thispagestyle{telosfirst}%
  \begingroup
  \setlength{\parskip}{0pt}%
  \noindent\rule{\linewidth}{1.6pt}\par
  \vspace{0.30em}
  \noindent
  \begin{minipage}[b]{0.66\linewidth}
    \raggedright{\LARGE\bfseries\TelosDocTitle}
  \end{minipage}%
  \hfill
  \begin{minipage}[b]{0.32\linewidth}
    \raggedleft{\small\itshape #1}
  \end{minipage}\par
  \vspace{0.18em}
  \noindent{\small\TelosDocSubtitle}\par
  \vspace{0.28em}
  \noindent\rule{\linewidth}{0.4pt}\par
  \endgroup
  \vspace{0.35em}
}

% ---------------------------------------------------------------------------
% Headings. Compact, small-caps, rule-anchored — a card has no room for the
% article class's default vertical air.
% ---------------------------------------------------------------------------
\titleformat{\section}{\normalfont\large\bfseries}{}{0pt}{}[\vspace{-0.55em}\rule{\linewidth}{0.4pt}]
\titlespacing*{\section}{0pt}{0.85em}{0.35em}
\titleformat{\subsection}{\normalfont\normalsize\bfseries}{}{0pt}{}
\titlespacing*{\subsection}{0pt}{0.6em}{0.2em}
\titleformat{\subsubsection}{\normalfont\small\bfseries\scshape}{}{0pt}{}
\titlespacing*{\subsubsection}{0pt}{0.45em}{0.12em}

% A short run-in heading for card interiors that must not consume a full line.
\newcommand{\telosrubric}[1]{\textbf{#1}\enspace}

% Cross-reference helper. Every Telos project is a set of loose printed sheets,
% so a reference names the sheet rather than a page number.
\newcommand{\sheetref}[1]{\textit{see the \textbf{#1} sheet}}

% Which decision, ADR or sheet governs the paragraph you are reading. A reader
% who disagrees with an instruction should be able to find the reasoning behind
% it rather than having to take it on trust.
\newcommand{\governedby}[1]{{\footnotesize\textsc{governed by #1}}}

% ---------------------------------------------------------------------------
% Quantities. siunitx is not in this TeX Live split, and the guides need only a
% fixed thin space and upright units, so define that directly rather than take
% a package dependency the documented install would not satisfy.
%   \qty{9}{kV}  ->  9 kV        \unitof{kV}  ->  kV
% ---------------------------------------------------------------------------
\newcommand{\unitof}[1]{\ensuremath{\mathrm{#1}}}
\newcommand{\qty}[2]{#1\,\unitof{#2}}
\newcommand{\numrange}[2]{#1\,--\,#2}
\newcommand{\qtyrange}[3]{#1\,--\,#2\,\unitof{#3}}
% Inches and feet appear constantly in the fishing guide; keep one spelling.
\newcommand{\inch}[1]{#1\,in}
\newcommand{\feet}[1]{#1\,ft}

% ---------------------------------------------------------------------------
% Cards and boxes.
% ---------------------------------------------------------------------------

% Titled card with a hairline frame. The default card for facts a reader looks
% up rather than reads through.
\newtcolorbox{teloscard}[1]{
  enhanced, breakable,
  colback=telospaper, colframe=telosink, colbacktitle=telospaper,
  coltitle=telosink,
  boxrule=0.5pt, sharp corners,
  left=6pt, right=6pt, top=4pt, bottom=4pt,
  toptitle=3pt, bottomtitle=3pt, titlerule=0.4pt,
  before skip=6pt, after skip=6pt,
  title={#1}, fonttitle=\small\bfseries
}

% Untitled left-ruled block: hierarchy without a redundant label.
\newtcolorbox{telosframe}{
  enhanced, breakable,
  colback=telospaper, frame hidden, boxrule=0pt, sharp corners,
  borderline west={1.3pt}{0pt}{telosink},
  left=8pt, right=2pt, top=3pt, bottom=3pt,
  before skip=6pt, after skip=6pt
}

% Filled emphasis panel — the "read this first" block. The 10% fill stays
% legible under a cheap laser toner curve and still reads as distinct.
\newtcolorbox{telospanel}[1]{
  enhanced, breakable,
  colback=telosfill, colframe=telosink, colbacktitle=telosfill,
  coltitle=telosink,
  boxrule=0.5pt, sharp corners,
  left=6pt, right=6pt, top=4pt, bottom=4pt,
  toptitle=3pt, bottomtitle=3pt, titlerule=0.4pt,
  before skip=6pt, after skip=6pt,
  title={#1}, fonttitle=\small\bfseries
}

% Hazard block. Deliberately the heaviest object on any page: a 1.4pt frame
% plus a solid title bar. Nothing else in Telos is allowed to look like this,
% so a reader skimming a lesson cannot mistake it for ordinary emphasis.
\newtcolorbox{teloshazard}[1]{
  enhanced, breakable,
  colback=telospaper, colframe=telosink,
  colbacktitle=telosink, coltitle=telospaper,
  boxrule=1.4pt, sharp corners,
  left=6pt, right=6pt, top=5pt, bottom=5pt,
  toptitle=3pt, bottomtitle=3pt,
  before skip=8pt, after skip=8pt,
  title={\MakeUppercase{#1}}, fonttitle=\small\bfseries
}

% A stop-and-check gate. Used before anything destructive, and before anything
% that would be expensive to discover later. Shared by every project that has
% a procedure someone could get hurt following carelessly.
\newtcolorbox{gate}[1]{
  enhanced, breakable,
  colback=telospaper, colframe=telosink,
  boxrule=1.0pt, sharp corners,
  left=6pt, right=6pt, top=4pt, bottom=4pt,
  toptitle=3pt, bottomtitle=3pt, titlerule=0.4pt,
  before skip=7pt, after skip=7pt,
  colbacktitle=telospaper, coltitle=telosink,
  title={\MakeUppercase{#1}}, fonttitle=\small\bfseries
}

% ---------------------------------------------------------------------------
% Rating dots. A five-step scale that prints identically in black-only: filled
% versus open circles, never a shaded bar.
% ---------------------------------------------------------------------------
\newcommand{\telosdot}{\tikz[baseline=-0.42ex]\fill[telosink] (0,0) circle (0.28ex);}
\newcommand{\telosopendot}{\tikz[baseline=-0.42ex]\draw[telosink,line width=0.28pt] (0,0) circle (0.28ex);}
\makeatletter
\newcounter{telos@rate}
\newcommand{\telosrating}[1]{%
  \begingroup
  \setcounter{telos@rate}{0}%
  \@whilenum\value{telos@rate}<#1\do{\telosdot\kern0.22em\stepcounter{telos@rate}}%
  \@whilenum\value{telos@rate}<5\do{\telosopendot\kern0.22em\stepcounter{telos@rate}}%
  \endgroup
}
\makeatother

% ---------------------------------------------------------------------------
% Tables.
%
% tabularx reads its body by scanning for a literal \end{tabularx}, so it can
% never be wrapped in a \newenvironment. Documents therefore open tabularx and
% longtable directly and take their house style from these column types and the
% booktabs rules. The one exception is the fact strip below, which is built on
% plain tabular precisely so it *can* be an environment.
% ---------------------------------------------------------------------------
\newcolumntype{L}{>{\raggedright\arraybackslash}X}
\newcolumntype{R}{>{\raggedleft\arraybackslash}X}
\newcolumntype{C}{>{\centering\arraybackslash}X}
\newcolumntype{B}{>{\bfseries\raggedright\arraybackslash}X}
\newcolumntype{N}{>{\raggedright\arraybackslash}p{0.16\linewidth}}

% Key/value fact strip. Two flush columns of short lookups; the label column is
% a fixed fraction so a stack of cards aligns across the whole guide.
\newenvironment{telosfacts}[1][0.24]
  {\begingroup\small
   \setlength{\parskip}{0pt}%
   \renewcommand{\arraystretch}{1.25}%
   \begin{tabular}{@{}>{\bfseries\raggedright\arraybackslash}p{\dimexpr#1\linewidth\relax}@{\hspace{0.7em}}>{\raggedright\arraybackslash}p{\dimexpr\linewidth-#1\linewidth-0.7em\relax}@{}}}
  {\end{tabular}\par\endgroup}
\newcommand{\telosfact}[2]{#1 & #2\\}

% ---------------------------------------------------------------------------
% Seasonal / diurnal activity bar.
%
% \telosseasonbar{4,3,2,1,1,2,3,3,4,5,4,3} draws a twelve-cell month strip
% where each value 0-5 is a fill height. Height, not gray value, carries the
% signal, so it is exact under any toner curve; the cell is outlined either way
% so an empty month is still visibly a month.
% ---------------------------------------------------------------------------
\newcommand{\telosseasonlabels}{J,F,M,A,M,J,J,A,S,O,N,D}
\newcommand{\telosseasonbar}[1]{%
  \begingroup
  \begin{tikzpicture}[x=1.25em,y=2.1ex,baseline=0pt]
    \foreach \v [count=\i from 0] in {#1} {
      \draw[telosrule,line width=0.3pt] (\i,0) rectangle (\i+0.86,5);
      \ifnum\v>0
        \fill[telosink] (\i,0) rectangle (\i+0.86,\v);
      \fi
    }
    \foreach \m [count=\i from 0] in \telosseasonlabels {
      \node[font=\tiny,anchor=north,inner sep=1.2pt] at (\i+0.43,-0.1) {\m};
    }
  \end{tikzpicture}%
  \endgroup
}

% Four-cell day-part strip: dawn, midday, dusk, night.
\newcommand{\telosdaylabels}{Dawn,Mid,Dusk,Night}
\newcommand{\telosdaybar}[1]{%
  \begingroup
  \begin{tikzpicture}[x=2.9em,y=2.1ex,baseline=0pt]
    \foreach \v [count=\i from 0] in {#1} {
      \draw[telosrule,line width=0.3pt] (\i,0) rectangle (\i+0.92,5);
      \ifnum\v>0
        \fill[telosink] (\i,0) rectangle (\i+0.92,\v);
      \fi
    }
    \foreach \m [count=\i from 0] in \telosdaylabels {
      \node[font=\tiny,anchor=north,inner sep=1.2pt] at (\i+0.46,-0.1) {\m};
    }
  \end{tikzpicture}%
  \endgroup
}

% A bar needs vertical room for its labels when it sits in a fact strip.
\newcommand{\telosbarrow}[1]{\rule{0pt}{4.2ex}#1\rule[-1.6ex]{0pt}{1.6ex}}

% ---------------------------------------------------------------------------
% Plates. A species or apparatus illustration with a caption rule beneath it.
% \telosplate{width fraction}{file}{caption}
% ---------------------------------------------------------------------------
\newcommand{\telosplate}[3]{%
  \begingroup
  \centering
  \includegraphics[width=#1\linewidth]{#2}\par
  \vspace{0.15em}
  \begin{minipage}{#1\linewidth}
    \rule{\linewidth}{0.4pt}\par
    \vspace{0.1em}
    {\footnotesize #3\par}
  \end{minipage}\par
  \endgroup
  \vspace{0.3em}
}

% TikZ drawing conventions shared by every rig, knot, and circuit diagram.
\tikzset{
  telos line/.style={draw=telosink,line width=0.7pt},
  telos hair/.style={draw=telosink,line width=0.35pt},
  telos heavy/.style={draw=telosink,line width=1.2pt},
  telos node/.style={draw=telosink,fill=telospaper,align=center,inner sep=4pt,font=\small},
  telos emphasis/.style={draw=telosink,line width=1.2pt,fill=telospaper,align=center,
    inner sep=4pt,font=\small\bfseries},
  telos label/.style={font=\scriptsize,inner sep=1.5pt},
  telos arrow/.style={-{Stealth[length=1.8mm]},draw=telosink,line width=0.6pt},
  telos dim/.style={|<->|,draw=telosink,line width=0.35pt},
  telos water/.style={draw=telosrule,line width=0.4pt,decorate,
    decoration={snake,amplitude=0.4mm,segment length=3mm}},
}

% ---------------------------------------------------------------------------
% Lesson furniture
% ---------------------------------------------------------------------------

% \lessonhead{number}{time}{who does what}
\newcommand{\lessonhead}[3]{%
  \par\noindent
  \begin{minipage}{\linewidth}
    \rule{\linewidth}{0.4pt}\par
    \vspace{0.3em}
    {\small\textbf{Lesson #1}\quad$\cdot$\quad #2\quad$\cdot$\quad #3\par}
    \vspace{0.25em}
    \rule{\linewidth}{0.4pt}
  \end{minipage}\par
  \vspace{0.5em}}

% The materials list. Two columns of short lines; a lesson that needs a long
% shopping list is a lesson that has not been thought through.
\newenvironment{materials}
  {\begin{telospanel}{What you need}\begin{multicols}{2}
   \begin{itemize}[leftmargin=1em,itemsep=0.05em,topsep=0pt]}
  {\end{itemize}\end{multicols}\end{telospanel}}

% The four rubrics every lesson carries, in the same order every time:
% what you are about to see, what to do, what it means, and what to ask.
\newcommand{\lessonsection}[1]{\section*{#1}\vspace{-0.35em}\rule{\linewidth}{0.4pt}\vspace{0.2em}}

% A prediction box. The point of the curriculum is that they commit to an
% answer before the demonstration, not after.
\newtcolorbox{predict}{
  enhanced, breakable,
  colback=telospaper, colframe=telosink,
  boxrule=0.5pt, sharp corners,
  left=6pt, right=6pt, top=4pt, bottom=4pt,
  before skip=6pt, after skip=6pt,
  borderline west={2.2pt}{0pt}{telosink},
  title={\small\bfseries Predict first}, fonttitle=\small\bfseries,
  colbacktitle=telospaper, coltitle=telosink,
  toptitle=3pt, bottomtitle=3pt, titlerule=0.4pt,
}

% ---------------------------------------------------------------------------
% Colophon. Rendered once, at the end of the last page, in the recovered footer
% area so it never creates a page of its own.
% ---------------------------------------------------------------------------
\newif\ifTelosColophonRendered
\newcommand{\TelosColophonExtra}{}
\newcommand{\TelosColophon}{%
  \ifTelosColophonRendered\else
  \global\TelosColophonRenderedtrue
  \par
  \enlargethispage{0.5in}%
  \vspace{0.3em}%
  \begingroup
  \setlength{\parskip}{0pt}%
  \begin{minipage}{\linewidth}
  \rule{\linewidth}{0.35pt}\par
  \vspace{0.3em}%
  \fontsize{7}{8.2}\selectfont
  \textbf{Telos.} \TelosProject{} — \TelosDocTitle. Revised \TelosRevision.
  \TelosColophonExtra{}
  Project-created text and diagrams are licensed \href{https://creativecommons.org/licenses/by/4.0/}{CC BY 4.0}.
  Incorporated public-domain artwork remains public domain and is credited on
  the plate it appears with. See \texttt{LICENSE} and \texttt{CREDITS.md} in the
  source. Nothing here is professional advice; verify current regulations and
  follow the stated safety procedure.
  \end{minipage}
  \endgroup
  \fi
}
\AtEndDocument{\TelosColophon}

% Shared macros for the homelab project.
%
% This project's documents are read in two very different situations: at a desk
% while designing, and standing in front of a broken machine at eleven at night
% with no network. The second case is the one that sets the style. Commands are
% set apart so they can be found by eye, every dangerous step has a visible
% verification after it, and nothing depends on being able to reach a web page.

% ---------------------------------------------------------------------------
% Decision-record entries, used by the generated digest.
% ---------------------------------------------------------------------------
\newcommand{\adrentry}[4]{%
  \par\vspace{0.7em}
  \needspace{4\baselineskip}%
  \noindent
  \begin{minipage}{\linewidth}
    \rule{\linewidth}{0.6pt}\par
    \vspace{0.25em}
    {\small\textbf{ADR #1}\quad #2\par}
    \vspace{0.1em}
    {\fontsize{7}{8.2}\selectfont\textsc{#3}\quad$\cdot$\quad #4\par}
    \vspace{0.2em}
    \rule{\linewidth}{0.3pt}
  \end{minipage}\par
  \vspace{0.3em}}

% ---------------------------------------------------------------------------
% Procedures.
%
% A rebuild manual is a sequence of commands and the observable evidence that
% each one worked. These environments keep that pairing structural rather than
% a matter of the author remembering.
% ---------------------------------------------------------------------------

% A block of shell to type. Deliberately not a listing package: plain verbatim
% in a ruled box prints identically everywhere and has no font dependency.
\newenvironment{shellblock}
  {\par\vspace{0.25em}\noindent
   \begin{minipage}{\linewidth}
   \rule{\linewidth}{0.3pt}\par\vspace{0.15em}
   \begingroup\footnotesize\ttfamily\obeylines\obeyspaces\parindent=0pt\parskip=0pt}
  {\endgroup\par\vspace{0.15em}\rule{\linewidth}{0.3pt}\end{minipage}\par\vspace{0.25em}}

% What you should see if the previous step worked. A step without one of these
% is a step you cannot verify, and in a rebuild manual that is a defect.
\newcommand{\expectthis}[1]{%
  \par\vspace{0.1em}
  \noindent\begin{minipage}{\linewidth}
  \hspace*{0.6em}\begin{minipage}{\dimexpr\linewidth-0.6em\relax}
  {\footnotesize\textbf{Expect:} #1\par}
  \end{minipage}\end{minipage}\par\vspace{0.3em}}

% \begin{gate} now lives in common/preamble.tex, shared with the other
% projects that document procedures.

% \governedby now lives in common/preamble.tex.

% An instance value that comes from the private overlay. In the published
% documents these render as visible placeholders, which is deliberate: a
% published manual must be complete and obviously generic, never look like it
% is describing somebody's actual network.
\newcommand{\instancevalue}[2]{\texttt{\textlangle #1\textrangle}%
  \footnote{Instance value. Supplied from \texttt{homelab/instance/} at build
  time; #2.}}
\newcommand{\placeholder}[1]{\texttt{\textlangle #1\textrangle}}

% ---------------------------------------------------------------------------
% Role and boundary diagrams.
% ---------------------------------------------------------------------------
\tikzset{
  role/.style={draw=telosink,line width=0.9pt,fill=telospaper,align=center,
    inner sep=5pt,font=\small,minimum width=2.6cm},
  roleoff/.style={draw=telosmute,line width=0.5pt,dash pattern=on 2.5pt off 1.8pt,
    fill=telospaper,align=center,inner sep=5pt,font=\small,minimum width=2.6cm},
  boundary/.style={draw=telosmute,line width=0.5pt,dash pattern=on 3pt off 2pt,
    rounded corners=3pt},
  flow/.style={-{Stealth[length=2mm]},draw=telosink,line width=0.7pt},
  hlabel/.style={font=\fontsize{6.4}{7.4}\selectfont},
}


\telosmeta{Homelab}{Convergence Design}
  {What day two looks like, and why installation is deliberately small}{2026-07-26}

\begin{document}
\telosmasthead[High-level design]

Installation does only what cannot be done later: partition, encrypt, write boot
artifacts, pin one interface, create one account. Everything else is a
repository, a commit and a converge run. A machine is not a thing that was set
up once; it is a thing that is continuously made to match what is written down.

\section{Where the line falls}
\begin{tabularx}{\linewidth}{@{}p{0.30\linewidth}L@{}}
\toprule
\textbf{Install time} & \textbf{Why it cannot wait}\\
\midrule
Partitioning and LUKS2 & Destructive, needs the authorization gate, and cannot be re-run on a machine in use. \governedby{ADR 0051, ADR 0058}\\
Boot artifacts and the UKI & The machine must boot before anything can converge it.\\
The managed interface pin & A MAC-pinned \texttt{.link} file has to exist before first boot decides whether to serve DHCP. \governedby{ADR 0050}\\
The manifest & Records what the machine \emph{is}. Everything afterwards reads it rather than being told. \governedby{ADR 0060}\\
\midrule
\textbf{Converge time} & \textbf{Why it belongs there}\\
\midrule
Packages, accounts, sshd & Changes constantly. A reinstall to add a package is a design failure.\\
dnsmasq, nginx, iPXE config & Generated from the manifest by the same code the installer uses, so a decision has one implementation.\\
Optional services & Off by default, enabled per instance, movable to another host. \governedby{ADR 0054}\\
Directory membership & Enabled only once the directory exists and answers. \governedby{ADR 0055}\\
\bottomrule
\end{tabularx}

\section{The manifest is the authority}
Convergence never tells a machine what it is. It asks.

\begin{center}
\begin{tikzpicture}[x=1cm,y=1cm]
  \node[role,minimum width=3.0cm] (repo) at (0,0) {repository\\playbooks + roles};
  \node[role,minimum width=3.0cm] (inv)  at (0,-1.7) {instance overlay\\gitignored};
  \node[telos emphasis,minimum width=3.0cm] (ans) at (4.2,-0.85) {ansible-playbook};
  \node[role,minimum width=3.0cm] (host) at (8.6,-0.85) {machine};
  \node[role,minimum width=3.0cm] (man)  at (13.0,-0.85) {manifest.json\\what it \emph{is}};
  \draw[flow] (repo) -- (ans);
  \draw[flow] (inv)  -- (ans);
  \draw[flow] (ans)  -- (host);
  \draw[flow] (host) -- (man);
  \draw[flow] (man.north) -- ++(0,0.8) -| (ans.north);
  \node[hlabel,anchor=south,align=center] at (8.6,0.55) {profile read back before anything is changed};
\end{tikzpicture}
\end{center}

\begin{telosfacts}[0.22]
\telosfact{The profile guard}{Each playbook reads \texttt{/etc/homelab/manifest.json} and refuses to proceed if the profile is not the one it converges. Converging a Workstation with the Controller playbook would start DHCP on it --- the one thing forbidden at every stage. \governedby{ADR 0008}}
\telosfact{Not inventory}{Inventory says how to \emph{reach} a machine. What the machine is comes from the machine. An inventory typo can therefore fail to connect, but cannot convert a Workstation into a DHCP server.}
\telosfact{Disposable is recorded}{A development-proof installation says so in its manifest, and convergence says so out loud. That is a recorded state, not a remembered one. \governedby{ADR 0043, ADR 0060}}
\end{telosfacts}

\section{One implementation of each decision}
The Controller's \texttt{dnsmasq}, \texttt{nginx} and iPXE configuration is not
an Ansible template. It is produced by the same generators the installer and the
bootstrap host call, invoked through a small bridge and installed verbatim.

\begin{telospanel}{Why this is worth the indirection}
A template is a second implementation. The generators refuse configurations that
violate a recorded decision --- a Controller address inside its own DHCP pool, a
default route advertised, \texttt{bind-dynamic} where \texttt{bind-interfaces} is
required --- and those refusals are unit-tested. A template that drifted would
still be valid \texttt{dnsmasq} syntax, would pass \texttt{dnsmasq -{}-test}, and
would quietly stop enforcing the decision it was written to enforce. The tests
assert byte equality between what convergence installs and what the generator
produces, so the drift cannot happen silently.
\end{telospanel}

\section{Getting back in when the directory is down}
The directory is a bootstrap dependency. While it is unavailable, only cached
logins and one local account work.

\begin{gate}{Convergence refuses to strand a machine}
Both the base role and the directory-join role assert that a break-glass
administrator key is configured, and stop if it is not. It is a dedicated key
pair used for nothing else: the daily-driver key is loaded in an agent all day
and forwarded to other hosts, so reusing it would protect the recovery account
with the credential most likely to be part of a bad day. Nothing in the
repository ever handles the private half. \governedby{ADR 0055, ADR 0063}
\end{gate}

The domain join itself is deliberately not automated. It requires
directory-administrator credentials, and automating it would mean storing those
somewhere a playbook can read unattended. The role detects that a machine is not
joined, stops, and prints the one command a person runs at a console.

\section{Public repository, private overlay}
\begin{telosfacts}[0.24]
\telosfact{In Git}{Playbooks, roles, generators, tests, decisions, and a template overlay in which every value is a visible placeholder.}
\telosfact{Not in Git}{Hostnames, addresses, interface MACs, disk serials, DHCP reservations and the break-glass public key. These live in a gitignored overlay that needs its own backup, because nothing else is backing it up. \governedby{ADR 0046}}
\telosfact{Not anywhere}{The LUKS2 passphrase, private keys, directory-administrator credentials and Kerberos keytabs. The overlay is gitignored, not encrypted, and sits in a working tree that gets copied around.}
\telosfact{Enforced, not intended}{The site build scans published sources for address literals and internal hostnames and fails closed. The template is tested to contain no key material. A tracked template that drifts from the variables the roles read is a test failure.}
\end{telosfacts}

\section{What is proven and what is not}
\begin{tabularx}{\linewidth}{@{}p{0.24\linewidth}p{0.12\linewidth}L@{}}
\toprule
\textbf{Stage} & \textbf{State} & \textbf{What it establishes}\\
\midrule
firmware & passing & A lab guest boots UEFI on its serial console, finds no boot disk, and attempts PXE over IPv4 on a segment with no route off it.\\
boot-chain & pending & The Controller answers DHCP, serves iPXE over TFTP, and the guest chainloads the installer over HTTP.\\
install & pending & The genuine installer is driven to completion over the guest's serial console, serial typed and all.\\
activation & pending & First boot refuses on a segment that already has DHCP, and starts cleanly on one that does not.\\
converge & pending & Ansible converges the installed Controller, and a second run reports no changes.\\
\bottomrule
\end{tabularx}

\begin{gate}{Pending is reported, never skipped}
A stage that cannot run yet prints what it is waiting for and the matrix does
not call itself green. A harness that silently skips what it cannot do converts
an unknown into a false assurance, which is worse than having no harness at all.
\end{gate}

\begin{telospanel}{A defect this design found before hardware did}
The lab's virtual disk reported no serial. The installer will not offer a disk
whose serial cannot be typed at the authorization prompt, so the acceptance
matrix could never have installed anything --- it would have failed with ``no
eligible disk'' on its first real run, and the obvious diagnosis would have been
that the installer was broken. The fix is one line of QEMU argv, but the reason
it was found is that the fixtures describe hardware \emph{shapes} rather than one
captured machine: NVMe, SATA and eMMC partition naming, absent serials,
removable media, wireless-only. A snapshot of one box would have missed it.
\end{telospanel}
\end{document}
